\documentclass[11pt]{article}
\usepackage[margin=1.1in]{geometry}
\usepackage{amsmath,amssymb,amsthm,booktabs}
\usepackage[hidelinks]{hyperref}
\usepackage[expansion=false]{microtype}
\setlength{\emergencystretch}{2em}
\newcommand{\ket}[1]{\lvert #1 \rangle}
\newcommand{\bra}[1]{\langle #1 \rvert}
\newcommand{\tr}{\operatorname{tr}}
\theoremstyle{definition}
\newtheorem{qprop}{Proposition}
\newtheorem{qlem}{Lemma}
\newtheorem{qrem}{Remark}
\title{Supplementary Material for\\ ``Database Schemas as Quantum Fields:\\ Relational and Document Models as a Classical Sector''}
\author{Arun Reungsilpkolkarn\\ School of Information Technology and Innovation, Bangkok University\\ \texttt{arun.r@bu.ac.th}}
\date{July 2026 \\[2pt] \small Code and data: \url{https://github.com/stevezhai139/qfdb-reproducibility} \\ Archived snapshot: \href{https://doi.org/10.5281/zenodo.21370700}{doi:10.5281/zenodo.21370700}}

\begin{document}
\maketitle

\begin{abstract}
\noindent This supplement contains what the six-page paper cannot: the division of labour between ``field'' and ``quantum'' (Appendix~A); a formal setup and \emph{full} proofs of the paper's propositions, including the coherence-extraction argument, the LHV bound, and the decoherence-functional additivity lemma (B); the complete algebra of the worked temporal example with the classical-disturbance derivation (C); the S1 power analysis (D); the full methods and results of the pre-registered stored-data negative control, including the settings chosen by the adversarial search and an instructive small-sample artifact that motivates the acceptance criterion (E); the LLM-respondent pilot in brief (F); and reproducibility details (G). The same document is distributed as the companion report in the repository, so the paper's references to ``the companion'' resolve here.
\end{abstract}
\appendix

\section{Why ``field'' and not just quantum mechanics}
The paper's three signatures (Bell/CHSH, order effects with the QQ invariant, interference) are properties of quantum \emph{theory} generally---two qubits suffice to violate CHSH---not of field theory specifically. The word \emph{field} is earned by a different part of the construction: data points are inserted and deleted, so the kinematics must support \emph{variable and superposable item count}, which is exactly what second quantization provides (Fock space, $a^\dagger/a$, the number operator as cardinality, fermionic single-valued vs.\ bosonic multi-valued modes). A fixed-dimensional quantum system cannot even represent a collection whose cardinality is in superposition (verified in \path{qfdb_H1_field_modes.py}, Part~4). The division of labour is deliberate: the Fock structure supplies the \emph{field}---it holds by construction and is not an empirical claim---while coherence and non-separability supply the \emph{quantum}, which is empirical, and the signatures test exactly that. Open on the field side, as stated in the paper's scope: the continuum/Gaussian lift, and genuine field dynamics beyond the update-channel semantics. A distinctly second-quantized empirical signature (exchange-statistics interference of indistinguishable excitations, in the style of Hong--Ou--Mandel) is a natural fourth test for a future version.

\section{Formal setup and full proofs}

\subsection{Setup}
Fix a finite \emph{structure basis} $\{\ket{s_k}\}_{k=1}^K$ of a Hilbert space $\mathcal H_\sigma$. An entity state is a density operator $\rho$ on $\mathcal H_\sigma$, decomposed as $\rho = D + C$ with $D=\sum_k p_k \ket{s_k}\bra{s_k}$, $p_k = \bra{s_k}\rho\ket{s_k}$, and $C=\rho-D$ (all off-diagonal terms). A \emph{query} is a projective measurement $\{P_a\}$ ($P_aP_{a'}=\delta_{aa'}P_a$, $\sum_a P_a=I$); a query outcome vector $\ket e$ denotes the rank-one case $P_e=\ket e\bra e$. Sequential measurement follows the standard rule: performing $\{P_a\}$ then $\{Q_b\}$ on $\rho$ yields the joint distribution
\begin{equation}\label{eq:seq}
p_{PQ}(a,b) \;=\; \tr\!\big(Q_b\,P_a\,\rho\,P_a\,Q_b\big),
\end{equation}
and the \emph{marginal of the second measurement} is $\Pr(b \text{ after } P)=\sum_a p_{PQ}(a,b)$.

\subsection{Proposition 1(a): the sector $\{C=0\}$ is closed under structure change}
\begin{proof}
A definite structure is $\ket{s_k}\bra{s_k}$; a heterogeneous or uncertain one is $\sum_k p_k\ket{s_k}\bra{s_k}$; both have $C=0$. Model schema evolution by a stochastic matrix $M$ ($M_{jk}\ge0$, $\sum_j M_{jk}=1$) acting on populations. The corresponding channel has Kraus operators $K_{jk}=\sqrt{M_{jk}}\,\ket{s_j}\bra{s_k}$:
\[
\mathcal E(\rho)=\sum_{j,k}M_{jk}\,\bra{s_k}\rho\ket{s_k}\,\ket{s_j}\bra{s_j},
\qquad
\sum_{j,k}K_{jk}^\dagger K_{jk}=\sum_k\Big(\sum_j M_{jk}\Big)\ket{s_k}\bra{s_k}=I,
\]
so $\mathcal E$ is completely positive and trace preserving. Its output is diagonal for \emph{every} input, and on diagonal inputs it acts as $p\mapsto Mp$. Lossless migrations (permutations $M$) are the special case $\mathcal E(\rho)=U_\pi\rho U_\pi^\dagger$ restricted to the diagonal. Hence no sequence of structure changes leaves $\{C=0\}$, and none creates coherence.
\end{proof}

\subsection{Proposition 1(b): which signature is carried by what}
\paragraph{Interference is carried by $C$, exactly.}
For an outcome $\ket e$,
\begin{equation}\label{eq:born}
P(e)=\bra e\rho\ket e=\underbrace{\sum_k p_k\,\lvert\langle e|s_k\rangle\rvert^2}_{\text{classical part}}\;+\;\underbrace{2\,\mathrm{Re}\!\!\sum_{k<l}C_{kl}\,\langle e|s_k\rangle\langle s_l|e\rangle}_{\text{interference part}} .
\end{equation}
If $C=0$ the interference part vanishes for every $e$. Conversely, suppose it vanishes for every $e$; fix $k<l$ and choose
\[
\ket{e_1}=\tfrac{1}{\sqrt2}(\ket{s_k}+\ket{s_l}),\qquad
\ket{e_2}=\tfrac{1}{\sqrt2}(\ket{s_k}+i\ket{s_l}).
\]
Both are supported on $\mathrm{span}\{s_k,s_l\}$, so only the $(k,l)$ term survives, giving interference parts $\mathrm{Re}\,C_{kl}$ and $-\mathrm{Im}\,C_{kl}$ respectively. Both vanish $\Rightarrow C_{kl}=0$ for all $k<l$, i.e.\ $C=0$. Hence the S3 signature is nonzero \emph{iff} $C\neq0$. \qed

\paragraph{The Bell signature is carried by non-separability, one-way.}
\begin{qlem}[Separable $\Rightarrow\lvert S\rvert\le2$]
Let $\rho=\sum_i p_i\,\rho_i^{A}\otimes\rho_i^{B}$ and let $\hat A,\hat A',\hat B,\hat B'$ be $\pm1$-valued observables. Then $\lvert S\rvert=\lvert E(A,B)+E(A,B')+E(A',B)-E(A',B')\rvert\le2$.
\end{qlem}
\begin{proof}
$E(X,Y)=\sum_i p_i\,a_i(X)\,b_i(Y)$ with $a_i(X)=\tr(\hat X\rho_i^A)\in[-1,1]$ and $b_i(Y)=\tr(\hat Y\rho_i^B)\in[-1,1]$. For each $i$,
\begin{align*}
\big\lvert a_i(A)\,[b_i(B)+b_i(B')]+a_i(A')\,[b_i(B)-b_i(B')]\big\rvert
&\le \lvert b_i(B)+b_i(B')\rvert+\lvert b_i(B)-b_i(B')\rvert\\
&= 2\max\big(\lvert b_i(B)\rvert,\lvert b_i(B')\rvert\big)\le2,
\end{align*}
and convexity over $p_i$ gives the bound. The converse is false (Remark~\ref{rem:werner}).
\end{proof}

\paragraph{Order effects are carried by non-commutation, not by $C$.}
\begin{qlem}[Commuting queries are order-invariant on every state]\label{lem:comm}
If $[P_a,Q_b]=0$ for all $a,b$, then for every $\rho$ the marginal of either measurement is unchanged by performing the other first.
\end{qlem}
\begin{proof}
Using \eqref{eq:seq}, cyclicity, and $P_aQ_bP_a=Q_bP_a^2=Q_bP_a$:
$\Pr(b \text{ after } P)=\sum_a \tr(Q_bP_a\rho P_a)=\sum_a\tr(P_aQ_bP_a\rho)=\sum_a \tr(Q_bP_a\rho)=\tr(Q_b\rho)$.
The symmetric claim follows by exchanging roles.
\end{proof}
\noindent Conversely, an order effect can occur on a \emph{coherence-free} state when the queries do not commute. Take $\rho=\ket0\bra0$ ($C=0$ in the stored basis) and $Q_\theta$ the projector onto $\cos\theta\ket0+\sin\theta\ket1$. Then
\[
p_{PQ}(0,\theta)=\cos^2\theta,
\qquad
p_{QP}(0,\theta)=\tr\big(P_0\,Q_\theta\rho Q_\theta\big)=\cos^2\theta\cdot\cos^2\theta=\cos^4\theta,
\]
which differ for $0<\theta<\pi/2$: the first measurement \emph{creates} the coherence (relative to the stored basis) that the second one feels. Face~A of the paper realises exactly this (Appendix~C). This is why the sector assigns order effects to the \emph{commuting} condition, interference to $C$, and Bell violations to non-separability---one signature per condition. \qed

\subsection{Proposition 2: data-model embedding, in full}
\begin{proof}
Build $\mathcal H$ by recursion on the type structure: an atomic field over domain $D$ is one mode $\mathbb C^{D\cup\{\bot\}}$ with orthonormal basis $\{\ket v\}_{v\in D\cup\{\bot\}}$; an ordered array of element type $\tau$ is the sequence space $\bigoplus_{k\ge0}\mathcal H_\tau^{\otimes k}$, and a set/multiset is its symmetric subspace; a nested document is the tensor product of its sub-fields. A record $t$ (relational tuple or document; heterogeneous collections $\bot$-padded) is the \emph{product} basis vector $\ket t=\bigotimes_f \ket{t_f}$. A relation or collection $r$ is
\[
\rho_r=\frac{1}{\lvert r\rvert}\sum_{t\in r}\ket t\bra t .
\]
(i)~\emph{Separability by definition:} $\rho_r$ is a convex combination of product states, which is the definition of separability across any factorisation of the field modes; no PPT/negativity criterion is invoked, so bound-entanglement counterexamples to ``PPT$\Rightarrow$separable'' (which exist from $2\times4$ and $3\times3$ upward) are irrelevant.
(ii)~\emph{Diagonality:} $\rho_r$ is diagonal in the record basis, so $C=0$; by Proposition~1(b) no interference signature exists.
(iii)~\emph{Commutation:} selection/membership projectors $\{\ket t\bra t\}$ are mutually orthogonal, hence commuting; by Lemma~\ref{lem:comm} the classical query family is order-invariant.
Codd's first normal form is the special case in which every field factor is a single mode; the document model enriches the \emph{space} (sequences, nesting, padding) while the state stays diagonal. Nothing in (i)--(iii) depends on which case holds, so both models occupy one diagonal, separable, commuting sector, and QFDB extends either by admitting off-diagonal, non-separable $\rho$ on the same $\mathcal H$.
\end{proof}
\begin{qrem}
In the paper's $2\times2$ example $\{(a,x),(b,y)\}$, quoting ``negativity $0$'' \emph{is} conclusive: PPT$\Leftrightarrow$separable holds in dimensions $2\times2$ and $2\times3$~\cite{horodecki1996}. The caveat in (i) applies only to higher dimensions.
\end{qrem}

\subsection{Proposition 3: temporal order-invariance, in full}
Histories formalism~\cite{griffiths1984}: for a family of histories $\{\alpha\}$ with chain operators $C_\alpha$ (time-ordered products of the projectors the history asserts, applied to the initial state), the decoherence functional is $D(\alpha,\beta)=\tr\big(C_\alpha\,\rho\,C_\beta^\dagger\big)$, with $P(\alpha)=D(\alpha,\alpha)$.

\begin{qlem}[Additivity $\Leftrightarrow$ weak consistency]\label{lem:add}
For $\alpha\neq\beta$ with coarse-graining $\gamma=\alpha\vee\beta$ (so $C_\gamma=C_\alpha+C_\beta$),
\[
P(\gamma)=D(\gamma,\gamma)=P(\alpha)+P(\beta)+2\,\mathrm{Re}\,D(\alpha,\beta).
\]
Hence the assignment $\alpha\mapsto P(\alpha)$ extends additively to all coarse-grainings of the family iff $\mathrm{Re}\,D(\alpha,\beta)=0$ for all $\alpha\neq\beta$.
\end{qlem}
\begin{proof} Expand $\tr\big((C_\alpha+C_\beta)\rho(C_\alpha+C_\beta)^\dagger\big)$ and use $D(\beta,\alpha)=\overline{D(\alpha,\beta)}$. \end{proof}

\begin{proof}[Proof of Proposition 3]
($\Leftarrow$)~If $\mathrm{Re}\,D(\alpha,\beta)=0$ pairwise, Lemma~\ref{lem:add} gives a single additive---hence classical---probability assignment on the whole Boolean algebra generated by the family; all query marginals are computed from that one joint distribution and are therefore independent of the order in which the queries are posed.
($\Rightarrow$)~If $\mathrm{Re}\,D(\alpha,\beta)\neq0$ for some pair, then by Lemma~\ref{lem:add} the ``probabilities'' fail additivity on $\alpha\vee\beta$, so no single classical joint distribution reproduces them, and order-dependence follows for generic queries distinguishing the pair (Appendix~C computes the dependence explicitly for Face~B).

\emph{Logging is sufficient:} a log is a \emph{record}. At the log time it applies the dephasing channel $\Delta(\sigma)=\sum_m P_m\sigma P_m$---equivalently, it writes the alternative into an ancilla, $\sigma\mapsto\sum_m P_m\sigma P_m\otimes\ket m\bra m_R$, and keeps the record. For histories $\alpha,\beta$ asserting different logged alternatives $m\neq m'$, the cross term of $D$ carries the record overlap $\langle m'|m\rangle_R=0$ (in channel form: $\Delta$ removes the $m\neq m'$ coherence before any later dynamics can act on it), so $D(\alpha,\beta)=0$ \emph{even though intervening evolution separates $P_m$ from $P_{m'}$ in the chain}---the naive step ``$P_mP_{m'}=0$'' alone would not suffice.

\emph{Logging is necessary exactly when the alternatives carry coherence:} if the pre-log state is already diagonal on the alternatives (a classical-sector state), the cross terms vanish with no logging---an ordinary unlogged classical database is order-safe. If the alternatives carry coherence with relative phase $\varphi$, then $\mathrm{Re}\,D(\alpha,\beta)\propto\cos\varphi\neq0$ for generic $\varphi$, and only logging (or dephasing) restores the condition. Hence: \emph{once histories carry coherence}, logging granularity is exactly the order-safety boundary.
\end{proof}

\subsection{The characterization is one-way, deliberately}
\begin{qrem}\label{rem:werner}
$\lvert S\rvert\le2$ does not imply separability: two-qubit Werner states are entangled in a visibility range where they satisfy the CHSH bound for \emph{all} settings (and in part of that range admit a full local-hidden-variable model)~\cite{werner1989}. Likewise, order-invariance of one realised query pair on one state does not imply the query algebra commutes; Lemma~\ref{lem:comm} quantifies over all states. The paper therefore states one-way guarantees: a state in the sector reproduces every classical value on these queries, and any signature beyond its classical bound certifies exit. The signatures are \emph{witnesses}, not characterisations of membership.
\end{qrem}

\section{Worked example: complete algebra}
Lifecycle: \emph{calf} $\to$ \emph{steer} $\to$ \emph{ox} (one male animal).

\paragraph{Face A (order effect on a stored record).}
Bloch convention: the stored axis is $\ket{\lnot\text{vacc}}=\ket0$; the mutant-susceptibility query projects onto $\ket{\theta}=\cos\tfrac\theta2\ket0+\sin\tfrac\theta2\ket1$ (Bloch angle $\theta$). Asked alone: $\Pr(\lnot\text{vacc})=1.000$. Asked after the mutant query, by \eqref{eq:seq}:
\[
\Pr(\lnot\text{vacc})=\cos^4\tfrac\theta2+\sin^4\tfrac\theta2
=1,\;0.875,\;0.750,\;0.625,\;0.500
\quad\text{at }\theta=0^\circ,30^\circ,45^\circ,60^\circ,90^\circ,
\]
an order effect up to $0.500$, while the Wang--Busemeyer QQ invariant~\cite{wangbusemeyer2013} is $0.000$ at every $\theta$ (\path{qfdb_order_effect_mutant.py}). A classical disturbance model tuned to the same order effect---asking B first redraws A fairly, B $\sim$ Bernoulli$(r)$---gives
\[
\mathrm{QQ}=\big[p(A_yB_y)+p(A_nB_n)\big]-\big[p(B_yA_y)+p(B_nA_n)\big]
=\big[r+0\big]-\big[\tfrac r2+\tfrac{1-r}2\big]=r-\tfrac12,
\]
i.e.\ $0.10,0.20,0.30$ at prevalence $r=0.6,0.7,0.8$: zero only at one fine-tuned base rate, whereas the quantum value is parameter-free.

\paragraph{Face B (unlogged history).}
Histories $\{$through-steer, not-through-steer$\}$ with amplitudes $a,b$ and relative phase $\varphi$. Unlogged, $\Pr(\text{ox now})=\lvert a+be^{i\varphi}\rvert^2\in[(\lvert a\rvert-\lvert b\rvert)^2,(\lvert a\rvert+\lvert b\rvert)^2]$; the observed swing $[0.000,0.444]$ forces $\lvert a\rvert=\lvert b\rvert$ and $4\lvert a\rvert^2=0.444$, so $\lvert a\rvert^2=1/9$. Logged, the decoherence functional dephases and $\Pr=\lvert a\rvert^2+\lvert b\rvert^2=2/9=0.222$, phase-blind; the off-diagonal magnitude of $D$ is $0.111$ unlogged and $0.000$ logged (\path{qfdb_consistent_histories.py})---the numerical instance of Lemma~\ref{lem:add}.

\paragraph{Complementarity (supporting check).}
$V^2+D^2\le1$~\cite{englert1996}: pinning records (which-structure distinguishability $D\to1$) sharpens the schema; leaving them indefinite keeps it field-like---a quantum reading of schema-on-write vs.\ schema-on-read. It holds for classical bits too, so it is a coherence check of the framework, not a non-classicality test.

\section{S1 power analysis}
Each $E(X,Y)$ is a mean of $n$ i.i.d.\ $\pm1$ products, so $\operatorname{Var}(E)\le(1-E^2)/n\le1/n$ and, with the four setting pairs assigned between-subjects (independent samples), $\operatorname{SE}(S)\le\sqrt{4/n}=2/\sqrt n$. To distinguish an Aerts--Sozzo-magnitude effect~\cite{aerts2011} ($S=2.42$, i.e.\ margin $0.42$ over the bound) at one-sided $\alpha=.05$ with power $.8$:
\begin{align*}
\frac{0.42}{2/\sqrt n}\;\ge\; z_{.95}+z_{.80}=1.645+0.842=2.487
\quad&\Longrightarrow\quad
\sqrt n\ge\frac{2\times2.487}{0.42}\approx11.8\\
&\Longrightarrow\quad
n\gtrsim140\ \text{per setting pair},
\end{align*}
the paper's ``on the order of $10^2$''. The bootstrap CI criterion is slightly more conservative than the normal approximation; a pre-registered interim look at half sample permits early stopping. Appendix~E's analysis~C shows empirically why small per-cell $n$ is dangerous. The elicitation design itself (parties, settings, coding, no-signaling gate, Level-2 hedge) is in the paper, \S8, and is not restated here.

\section{Pre-registered stored-data negative control: full methods and results}
Pre-registration frozen before evaluation (\path{s1_stored/PREREGISTRATION.md}, sha256 prefix \texttt{72cf065266be0124}). Prediction: stored two-store correlations are common-cause, hence $S\le2$ everywhere including under adversarial setting search; for per-side-deterministic responses this is a theorem (the shared item is the latent variable of the Lemma in B.3; cf.~Fine~\cite{fine1982}), so the run validates the instrument end-to-end and measures the margin.

\paragraph{Data and responses.} Magellan/ditto entity-matching benchmarks; items $=$ ground-truth matched pairs; rules built per side on TRAIN only (numeric-median, length-median, top-8 title tokens); evaluation on held-out VALID$\cup$TEST; missing value $\to-1$ (frozen). Locality by construction: each side's $\pm1$ response is a function of that side's record and setting only.

\paragraph{Analyses.} (A)~natural settings on the same held-out items; (B)~\emph{adversarial}: exhaustive search over ordered rule pairs maximising $\lvert S\rvert$ on TRAIN, argmax evaluated once held-out with a $10^4$-resample bootstrap CI; (C)~between-subjects dry run: held-out items split into four disjoint groups (seed 42), one setting pair each, with no-signaling deltas. Self-tests pass before data: Bell $2.8284$; engineered classical correlations attain $S=2.000$ exactly (the instrument is sharp at the bound); product state $\approx0$.

\begin{table}[h]\centering\small
\caption{Held-out results. Analysis B is the bound test: no CI lies above $2$. The entangled control in the same pipeline reaches $2.8284$.}
\begin{tabular}{lrrrrcc}
\toprule
dataset & $n$ & $S_{\text{nat}}$ & $S_{\text{adv}}^{\text{train}}$ & $S_{\text{adv}}^{\text{held-out}}$ & $95\%$ CI & NS $\Delta$ (C)\\
\midrule
Amazon--Google & 468 & $+1.026$ & $+1.966$ & $+1.974$ & $[+1.940,+2.000]$ & $0.154$\\
DBLP--ACM & 888 & $-0.189$ & $+2.000$ & $+1.995$ & $[+1.986,+2.000]$ & $0.189$\\
DBLP--GoogleScholar & 2140 & $-0.456$ & $+1.988$ & $+1.983$ & $[+1.972,+1.993]$ & $0.037$\\
Fodors--Zagats & 44 & $-0.364$ & $+2.000$ & $+2.000$ & $[+2.000,+2.000]$ & $0.364$\\
Walmart--Amazon & 386 & $-0.166$ & $+1.979$ & $+1.938$ & $[+1.886,+1.979]$ & $0.103$\\
iTunes--Amazon & 54 & $+0.815$ & $+2.000$ & $+2.000$ & $[+2.000,+2.000]$ & $0.319$\\
\bottomrule
\end{tabular}
\end{table}

\paragraph{What the adversary chose.} The argmax settings $(A,A';\,B,B')$ are instructive:

\begin{center}\small
\begin{tabular}{ll}
Fodors--Zagats & (\texttt{tok:cafe}, \texttt{tok:room};\ \texttt{tok:room}, \texttt{tok:cafe})\\
DBLP--GoogleScholar & (\texttt{tok:mining}, \texttt{tok:queries};\ \texttt{tok:queries}, \texttt{tok:mining})\\
iTunes--Amazon & (\texttt{tok:love}, \texttt{tok:feat};\ \texttt{tok:feat}, \texttt{tok:love})\\
Amazon--Google & (\texttt{tok:adobe}, \texttt{tok:home};\ \texttt{tok:punch}, \texttt{tok:adobe})\\
\end{tabular}
\end{center}

\noindent The optimizer converges on \emph{mirrored} rules across the two sides: it rediscovers the textbook extremal classical strategy of near-perfectly correlated responses, which saturates but cannot exceed $2$ (B.3). Real integration data thus \emph{crowds} the classical ceiling exactly as the embedding requires.

\paragraph{An instructive small-sample artifact.} In analysis C the four $E$'s are estimated on \emph{disjoint} small groups ($n/4$ items each). On the small datasets the naive point estimate then exceeds $2$ spuriously---Fodors--Zagats $S_{\text{C}}=2.364$, iTunes--Amazon $2.330$, DBLP--ACM $2.090$---accompanied by exactly the tell the protocol looks for: large no-signaling deltas ($0.36$, $0.32$, $0.19$), i.e.\ marginals that shift with the distant setting because different items sit in different cells. The pre-registered criterion (CI clearance \emph{and} no-signaling pass) fires on none of them. This is the empirical justification for both halves of the acceptance criterion and for the sample sizes of Appendix~D; a human study with a dozen respondents per cell could manufacture ``violations'' the same way.

\paragraph{Disclosure.} Beer ($\approx20$ held-out matches) is excluded by a $\ge30$-item practicality rule adopted at analysis time and disclosed here (not part of the frozen pre-registration); all other pairs are reported without selection. Reproduce with \path{s1_stored/s1_stored_negative_control.py} (fetch commands in the pre-registration).

\section{LLM-respondent pilot (pointer)}
\path{s1_elicitation/llm_harness/} implements the paper's S1 protocol with a language model as synthetic respondent (one call $=$ one respondent; one setting pair per call; neutral wording; counterbalanced options; temperature grid; same acceptance criterion). Its roles are instrument validation and an effect-size prior for the human study; a violation would certify non-classical structure in the \emph{model's} judgment statistics only. Method discipline and caveats are documented in \path{README_llm.md}; a \texttt{--dry-run} classical judge exercises the analysis path without any API.

\section{Reproducibility}
All results reproduce digit-for-digit (including bootstrap CIs; fixed PCG64 seed) on two independent machines: (a)~MacBook Pro (Apple M4, 24\,GB RAM, 512\,GB internal SSD $+$ 1\,TB Thunderbolt external SSD), macOS Tahoe~26.x, NumPy~2.2.6; (b)~Debian Linux aarch64, Python~3.10, NumPy~2.2.6. No GPU or quantum hardware; the full negative-control run completes in under a minute. Scripts are NumPy-only (the LLM harness is stdlib-only): \path{qfdb_signatures_verify.py}, \path{qfdb_order_effect_mutant.py}, \path{qfdb_consistent_histories.py}, \path{qfdb_defs_verify.py}, \path{qfdb_commoncause_bound.py}, \path{qfdb_H1_field_modes.py} reproduce every computed number in the paper; \path{s2_calibration/} reproduces the 72-survey calibration~\cite{wang2014} (order-effect mean $0.087$, range $0.03$--$0.58$; QQ mean \mbox{$+0.001\pm0.036$}, $t(71)=0.25$, $p=0.80$); \path{s1_stored/} reproduces Appendix~E.

\begin{thebibliography}{9}
\bibitem{werner1989} R.~F. Werner. Quantum states with Einstein--Podolsky--Rosen correlations admitting a hidden-variable model. \emph{Physical Review A} 40(8):4277--4281, 1989.
\bibitem{horodecki1996} M.~Horodecki, P.~Horodecki, and R.~Horodecki. Separability of mixed states: necessary and sufficient conditions. \emph{Physics Letters A} 223(1--2):1--8, 1996.
\bibitem{fine1982} A.~Fine. Hidden variables, joint probability, and the Bell inequalities. \emph{Physical Review Letters} 48(5):291--295, 1982.
\bibitem{griffiths1984} R.~B. Griffiths. Consistent histories and the interpretation of quantum mechanics. \emph{Journal of Statistical Physics} 36(1--2):219--272, 1984.
\bibitem{englert1996} B.-G. Englert. Fringe visibility and which-way information: an inequality. \emph{Physical Review Letters} 77(11):2154--2157, 1996.
\bibitem{wangbusemeyer2013} Z.~Wang and J.~R. Busemeyer. A quantum question order model supported by empirical tests of an a priori and precise prediction. \emph{Topics in Cognitive Science} 5(4):689--710, 2013.
\bibitem{wang2014} Z.~Wang, T.~Solloway, R.~M. Shiffrin, and J.~R. Busemeyer. Context effects produced by question orders reveal quantum nature of human judgments. \emph{PNAS} 111(26):9431--9436, 2014.
\bibitem{aerts2011} D.~Aerts and S.~Sozzo. Quantum structure in cognition: Why and how concepts are entangled. In \emph{Quantum Interaction (QI 2011)}, LNCS 7052, pages 116--127. Springer, 2011.
\end{thebibliography}

\end{document}
